\begin{abcliste}
\abc Between the plates the field lines are parallel and evenly spaced: the field is uniform, as strong near the plates as in the middle. At the edge the lines bulge out and spread apart: the field is weaker there. Outside the capacitor the fields of the two plates cancel: there is (almost) no field.
\begin{align}
 \result{E_P = E_Q = E_R > E_S > E_T \approx 0}
\end{align}

\abc The field lines leave the wire at right angles and spread out like the spokes of a wheel, but only in the plane across the wire, not in all directions as round a point charge. At twice the distance the same lines are spread over a circle with twice the circumference, so they are half as dense:
\begin{align}
 E = \frac{\lambda}{2\pi\varepsilon_0\,r} \qquad\Rightarrow\qquad \frac{E(\SI{4}{\cm})}{E(\SI{2}{\cm})} = \frac{2}{4} = \result{\frac{1}{2}}
\end{align}
(A quarter would be the answer for a point charge, whose field falls off as $1/r^2$.)
\end{abcliste}
